\bookchapter{第五十六章：打磨同發布——測試、除錯同上架 Google Play}

\chapterintro{App 做咗出嚟，跟住點呢？呢一章回答「由跑得郁」到「可以交付」嘅三件事：點樣高效咁除錯搵 bug、點樣寫測試防走樣、點樣將 App 攞上 Google Play 畀全世界下載。上架唔只係儀式感——佢係將一件作品變成產品嘅分水嶺。}

\booksection{一、除錯三板斧：Logcat、斷點、UI 檢查器}

\begin{center}
\begin{tabularx}{\linewidth}{|X|X|X|}
\hline
\textbf{工具} & \textbf{做乜嘢} & \textbf{點用} \tabularnewline
\hline
Logcat & 睇系統/自己打嘅日誌 & 代碼入面 Log.d("TAG","值=\$x")，喺面板搜 TAG 過濾 \tabularnewline
斷點除錯 & 令程式停喺某一行，睇變數值 & 喺行號旁邊㩒一下出紅點，用 Debug 模式行，逐步咁行 \tabularnewline
Layout Inspector & 睇運行中嘅界面結構 & 選單 Tools Layout Inspector，實時睇組件樹 \tabularnewline
\hline
\end{tabularx}
\end{center}

\begin{center}\small\color{gray}除錯三板斧\end{center}

\begin{tcbinfo}{小貼士：報錯點樣睇？}
崩潰之後 Logcat 會有紅色堆疊（Stacktrace），由下往上搵第一行以自己包名開頭嘅（例如 com.example.yourapp），嗰行基本上就係案發現場。先睇「Exception 類型」（空指標？類型轉換？），再睇「at ... 你嘅檔案:行號」。呢一招可以解決 80\% 嘅崩潰。

\end{tcbinfo}

\booksection{二、測試：畀代碼買保險}

手動測試靠唔住——改咗呢個功能，另一個靜靜雞壞咗。測試就係「後悔藥」兼「防彈衣」。安卓測試分三層，新手先讀透第一層：

\begin{center}
\begin{tabularx}{\linewidth}{|X|X|X|X|}
\hline
\textbf{層級} & \textbf{測乜嘢} & \textbf{擺喺邊} & \textbf{速度} \tabularnewline
\hline
單元測試 & 純邏輯（計算、校驗） & app/src/test/ & 秒級 \tabularnewline
儀器測試 & 真實設備上嘅界面行為 & app/src/androidTest/ & 分鐘級 \tabularnewline
手動測試 & 真人感受、體驗 & —— & 永遠都需要 \tabularnewline
\hline
\end{tabularx}
\end{center}

\begin{center}\small\color{gray}測試三層\end{center}

\begin{lstlisting}[language=Kotlin]
// 把"金額字串轉數字"呢種邏輯抽成純函數
fun parseAmount(input: String): Double? {
    val v = input.trim().toDoubleOrNull()
    return if (v != null && v > 0) v else null
}

// 測試佢（JUnit4）
@Test
fun emptyInputReturnsNull() {
    assertEquals(null, parseAmount(""))
}
@Test
fun negativeReturnsNull() {
    assertEquals(null, parseAmount("-5"))
}
@Test
fun normalValueParses() {
    assertEquals(12.5, parseAmount("12.5"))
}
\end{lstlisting}

記賬本入面嘅金額校驗、統計計算，都應該有單元測試罩住。行測試：右鍵 test 目錄 Run。全部綠色過關，就可以放心改代碼。

\booksection{三、簽名打包：畀 App 蓋上「印章」}

安卓要求所有安裝包必須簽名——就好似蓋公章咁，證明「呢個包係我出嘅，冇畀人改過」。除錯嗰陣 Android Studio 會自動用 debug 簽名；發布就要用自己嘅 release 簽名。流程係：

\begin{enumerate}
  \item 選單 Build Generate Signed Bundle / APK；
  \item 建立 keystore 檔案（密鑰庫）——設定密碼同埋妥善保管；
  \item 揀 release + AAB 格式，構建完成。
\end{enumerate}

\textbf{避坑提醒：keystore 唔見咗 = App 唔見咗}\par

上架之後更新 App 必須用同一個簽名，否則商店會拒絕更新。keystore 檔案同密碼一旦唔見咗，你就只可以另開一個新應用重新上架——舊用戶全部接唔返。一定要將 keystore 同密碼備份去安全嘅地方（雲盤加密 + 記事本），呢個係安卓開發者最昂貴嘅教訓之一。

\booksection{四、上架 Google Play 全流程}

註冊開發者賬號

準備商店素材

上傳 AAB

填數據安全表

審覈（數小時\textasciitilde{}數天）

上架

\begin{center}\small\color{gray}圖 6：Google Play 上架六步走\end{center}

\begin{itemize}
  \item 註冊：一次性 25 美元註冊費，需要一個用得嘅 Google 賬號；
  \item 素材：應用圖標（512×512）、截圖（最少 2 張）、功能描述、私隱政策連結（有收集數據就必填）；
  \item 數據安全表：如實申報收集咗邊啲數據——唔聯網嘅記賬本可以「零收集」；
  \item 審核：新賬號首次上架好多時會經人工審核，幾個鐘到幾日唔等，被打回頭就照意見改。
\end{itemize}

國內用戶注意：Google Play 需要特殊網絡環境，而且中國大陸地區分發通常行國內商店（華為、小米、應用寶等）。上架國內商店嘅流程類似，但各商店有自己嘅審核口味——同一個 AAB/APK 可以多間店一齊分發。作為練手，先行通一間就已經贏咗。

\booksection{五、商店素材點樣做：第一印象就係下載量}

商店頁係用戶「行店」嗰陣見到嘅櫥窗。三樣素材決定第一印象：

\begin{itemize}
  \item 圖標：512×512，唔好攞截圖求其算數。冇設計能力就用 AI 工具生成一個簡潔幾何圖形，喺手機桌面縮細之後都仲可以一眼認得出先算合格；
  \item 截圖：最少 2 張，最好 4-6 張，展示最核心嘅功能場景；每張配一句講人話嘅文案（「一秒記一筆」好過「高效嘅記賬解決方案」十倍）；
  \item 描述：第一句講清楚「呢個係乜嘢 + 畀邊個用」，唔好亂堆「最專業、最勁」——商店搜索權重主要睇關鍵詞同下載轉化，真誠好過誇張。
\end{itemize}

\begin{tcbinfo}{小貼士：A/B 睇吓同行}
上架之前行一圈商店入面同類 App 嘅前十名，截圖嚟對比：佢哋啲圖標係咩風格？截圖講咩賣點？標題點寫？模仿唔係抄襲——學套路、換自己做嘅內容，呢個係所有產品冷啟動最快嘅一課。

\end{tcbinfo}

\booksection{五、性能優化同崩潰上報小清單}

App 上架之前，呢兩件事值得各花半個鐘：

\booksubsection{性能自查四條}

\begin{enumerate}
  \item 列表用 LazyColumn：長列表千祈唔好用 Column 硬排——懶加載係安卓性能嘅第一課；
  \item 圖片唔好硬塞原圖：用 Coil 圖片庫自動縮放同緩存，一張 5MB 嘅原圖足以卡爆低端機；
  \item 避免喺主線程做重活：數據庫、網絡、大檔案都擺喺協程度做；
  \item 用 Profiler 度一度：Android Studio 自帶嘅性能分析器睇到內存、CPU、耗電——先量度，再優化，唔好齋靠感覺。
\end{enumerate}

\booksubsection{崩潰上報}

你測試嗰陣冇崩，唔代表用戶部手機唔崩。上架之前接入 Firebase Crashlytics（谷歌官方，免費）：幾行配置，用戶崩潰嗰陣自動上報堆疊，你喺後台就會見到「邊個型號、邊個版本、崩喺邊一行」。冇崩潰監控就上架，等於瞇埋眼揸車。

\begin{tcbpractice}{嚟動手：將記賬本「交付」一次}
(1) 幫 parseAmount 呢啲邏輯補 3 個單元測試；(2) 用 Logcat 親手捉一個自己整出嚟嘅崩潰然後修好佢；(3) 生成一個 release 簽名嘅 APK，發畀朋友安裝（呢個先係「作品」嘅儀式感）；(4) 有餘力就註冊開發者賬號，將記賬本真正提交上架。前三條今晚就做得完。

\end{tcbpractice}

\begin{tcbsummary}{本章小結}
\begin{itemize}
  \item 除錯三板斧：Logcat 睇日誌、斷點睇變數、Layout Inspector 睇界面
  \item 讀崩潰：由堆疊下面往上搵自己包名嘅第一行
  \item 測試三層：單元（秒級，先學呢個）/儀器/手動
  \item 發布必須簽名：keystore 同密碼備份好，唔見咗個 App 就等於「換咗人」喇
  \item 上架六步：註冊 素材 上傳 AAB 數據安全表 審核 上架
\end{itemize}

\end{tcbsummary}
